% Structure follows 6.pdf; all six sections use audited implementation/evidence.
\chapter{Workload-Level Refinement-Budget Calibration}
\label{ch:refinement}

Chapters 4 and 5 establish a compressed index and its GPU execution.
This chapter addresses RQ3: how much refinement work a workload needs,
and whether selecting that budget repays its calibration cost. The
policy compares unlabelled query outputs with a larger-budget reference
and reuses one selected multiplier across batches. The method is analysed
through its acceptance rule, conditional selection theorem, experimental
controls and amortisation model. Chapter 7 reports the measured outcomes.

\section{Why Fixed Refinement Budgets Can Be Wasteful}
\label{sec:6-1}
JHQ screens routed candidates with primary codes before reading residuals
for a smaller survivor set \cite{jhq}. For a query $q$, let $U(q)$ be
the candidates in its selected IVF lists. The refinement multiplier sets
\begin{equation}
 c(\alpha)=\max\{k,\lceil\alpha k\rceil\},\qquad
 C_\alpha(q)=\operatorname{Top}^{\rm primary}_{\min\{c(\alpha),|U(q)|\}}U(q).
 \label{eq:6-budget}
\end{equation}
Here $C_\alpha$ denotes valid survivors. The implementation reserves
$c(\alpha)$ slots and can pad an underfull pool with sentinels. The
analysis assumes at least $k$ valid candidates. Raising $\alpha$ expands
refinement capacity within the visited lists; raising $n_{\rm probe}$
changes the visited universe. Neither operation recovers information
lost by the other automatically.

A large fixed $\alpha$ may repeatedly refine candidates that do not
change the final top-$k$. A small value can exclude a useful neighbour
before residual scoring. Which tradeoff is worthwhile depends on the
query population, primary reconstruction, list coverage and refined-score
ordering. One global default therefore need not suit every dataset--probe
configuration. The evaluated grid spans $c=20$ to $2{,}000$ at $k=10$.
Changing it reuses the same trained index and stored residuals.
Lowering the multiplier reduces residual rows accessed, rather than
releasing the resident residual representation. Survivor buffers grow
when capacity is insufficient and retain that capacity afterward, so a
smaller budget is not a proportional reduction in allocated memory.

\begin{figure}[tbp]
\centering
% Logical work boundary, not a measured latency or pruning diagram.
\begin{tikzpicture}[
 font=\footnotesize\linespread{1}\selectfont,
 box/.style={draw=ThesisNavy!65,rounded corners=2pt,align=center,
 text width=2.15cm,minimum height=1.5cm,inner sep=3pt},
 arr/.style={-{Latex[length=2mm]},draw=ThesisNavy!80,thick}]
\node[box] (route) at (0,0) {Probe IVF lists\\Fixed $n_{\rm probe}$};
\node[box] (primary) at (3,0) {Score $U(q)$\\Primary codes};
\node[box,fill=ThesisNavy!7] (select) at (6,0)
 {Primary top-$c$\\Keep $C_\alpha(q)$};
\node[box,fill=ThesisNavy!7] (refine) at (9,0)
 {Refine $c(\alpha)$ slots\\Residual rows};
\node[box,fill=ThesisNavy!7] (final) at (12,0)
 {Final top-$k$\\Return $k$ IDs};
\draw[arr] (route) -- (primary);
\draw[arr] (primary) -- (select);
\draw[arr] (select) -- (refine);
\draw[arr] (refine) -- (final);
\node[align=center,text width=5.1cm] at (1.5,1.4)
 {Fixed routed universe and number of primary scores};
\node[align=center,text width=7.6cm] at (9,1.4)
 {Budget-dependent selection, residual scoring and final selection};
\draw[densely dashed,ThesisNavy!45] (4.5,-0.9) -- (4.5,1.1);
\node[align=center,text width=13.8cm] at (6,-1.45)
 {$c(\alpha)=\max\{k,\lceil\alpha k\rceil\}$ slots;
 valid survivors are capped by $|U(q)|$.\\
 The stored model and residual codes are reused across budgets.};
\end{tikzpicture}

\caption[Work affected by the refinement budget]{The refinement budget changes work after primary-score
production. With a fixed query, index and probe count, it leaves the
routed candidate universe unchanged. The highlighted stages depend on
survivor capacity; this is a logical work decomposition, not a measured
latency breakdown. The scan kernel also performs budget-dependent
admission and compaction. Padding applies when fewer than $c$ candidates exist.}
\label{fig:6-budget-scope}
\end{figure}

An explanatory latency decomposition is
\begin{equation}
 T(\alpha)=T_{\rm base}+T_{\rm select}(c(\alpha))+
 T_{\rm refine}(c(\alpha))+T_{\rm final}(c(\alpha)).
 \label{eq:6-cost-model}
\end{equation}
For a fixed batch and route, $T_{\rm base}$ includes transfer, rotation,
routing and primary-score production; selection, residual scoring and
final selection depend on survivor capacity (\cref{fig:6-budget-scope}).
Buffer capacities, occupancy and resource thresholds can change in steps,
so this work model does not predict a monotone runtime. The measured
small-budget plateau is evaluated in \cref{sec:7-budget-trace}; the policy
minimises accepted budget rather than a noisy timing rank.


Each calibration trial still runs rotation, probing and primary scoring
before refinement. Reusing the index avoids rebuilding or retraining;
it does not cache every intermediate query result across budget trials.
The cost of calibration must therefore count full searches, not merely
the difference between two residual-scoring workloads.

The policy selects one budget for an index, routing configuration, $k$
and query workload. It uses neither a trained per-query predictor nor
an online choice for each arriving query. Scan-group scheduling remains
the separate mechanism of \cref{sec:5-5}; its execution path must still
be considered when interpreting budget calibration.

\paragraph{Position among adaptive ANN methods.}
The multiplier has the same candidate-pool role as FAISS
\texttt{IndexRefine::k\_factor} and cuVS's refinement ratio
\cite{faiss_refine,cuvs_refinement}. FAISS auto-tuning explores
speed--accuracy operating points using a ground-truth criterion
\cite{faiss_autotune}. Li et al. learn per-query termination conditions
from search features \cite{li_termination}; DARTH predicts recall during
execution to stop at a requested target \cite{darth}.
Here the decision calibrates one unlabelled workload budget for reuse,
using reference-output coverage rather than a trained recall predictor.
The final scorer remains a compressed reconstruction, whereas the cuVS
refinement example reranks source vectors. These distinctions separate
the work parameter, its selection mechanism and the scored representation.


\section{GT-Free Sampled Output Agreement}
\label{sec:6-2}
Let $\mathcal R_\alpha(q)$ be the set of $k$ valid, unique database IDs
returned at budget $\alpha$ on a fixed index. Define
$\mathcal R_{\rm ref}(q)=\mathcal R_{\alpha_{\rm ref}}(q)$. The descending grid is
\begin{equation}
 A=(200,100,64,32,16,8,4,2),\qquad \alpha_{\rm ref}=A[0]=200.
 \label{eq:6-grid}
\end{equation}
For $S$ calibration queries $\mathcal Q_S$, define
\begin{equation}
 D_S(\alpha)=\sum_{q\in\mathcal Q_S}
 \left(k-|\mathcal R_\alpha(q)\cap\mathcal R_{\rm ref}(q)|\right),
 \qquad D_S(\alpha)\leq\epsilon.
 \label{eq:6-agreement}
\end{equation}
The evaluated $\epsilon=1$ permits one missing reference ID across the
whole sample. It ignores output order and distance differences; replacing
one ID counts once, although the symmetric difference contains two IDs.
The finite-grid objective is
\begin{equation}
 \widehat\alpha=\min\{\alpha\in A:D_S(\alpha)\leq\epsilon\}.
 \label{eq:6-selection-objective}
\end{equation}
The reference supplies a nonempty fallback. The objective minimises an
accepted multiplier; \cref{sec:6-3} states when bisection solves it exactly.

At $k=10$, the normalised allowance $\epsilon/(Sk)$ is 0.003125,
0.0015625 and 0.00078125 for $S=32,64,128$, and 0.0002 at $S=500$.
Increasing $S$ therefore changes both coverage and effective strictness.
Calibration and held-out evaluation use disjoint queries; subsamples
reuse a finite calibration pool (\cref{sec:7-5}).

\begin{figure}[htbp]
\centering
% Native schematic of the implemented decision and its offline evaluation.
\begin{tikzpicture}[
 font=\footnotesize\linespread{1}\selectfont,
 box/.style={draw=ThesisNavy!65,rounded corners=2pt,align=center,
 text width=2.8cm,minimum height=1.7cm,inner sep=4pt},
 arr/.style={-{Latex[length=2mm]},draw=ThesisNavy!80,thick}]
\node[box] (sample) at (0,0) {Calibration pool\\Even indices\\Sample $S$ queries};
\node[box,fill=ThesisNavy!7] (ref) at (3.6,0) {Reference search\\$\alpha_{\rm ref}=200$\\Top-$k$ IDs};
\node[box,fill=ThesisNavy!7] (test) at (7.2,0) {Bisection\\$D_S(\alpha)\leq1$\\Shared sample};
\node[box] (choice) at (10.8,0) {Chosen $\alpha$\\Caller applies it\\For batch reuse};
\node[box] (held) at (0,-3) {Held-out pool\\Odd query indices\\GT for evaluation};
\node[box,text width=6.4cm] (eval) at (7.2,-3)
 {Evaluate the selected arm\\Look up held-out recall and QPS\\from previously timed grid arms};
\draw[arr] (sample) -- (ref);
\draw[arr] (ref) -- (test);
\draw[arr] (test) -- (choice);
\draw[arr] (held) -- (eval);
\draw[arr,dashed] (choice.south) -- (10.8,-3) -- (eval.east);
\end{tikzpicture}

\caption{Workload-level calibration and its evaluation. Budget decisions
use reference IDs; GT assesses held-out quality. Selected-arm results
are looked up from timed fixed arms. Implementation boundaries are in
\cref{sec:a-budget-contract}.}
\label{fig:6-calibration}
\end{figure}

\subsection{What agreement controls}

Under \Cref{prop:6-monotonicity}, disagreement counts reference answers
excluded by primary screening:
\begin{equation}
 D_S(\alpha)=\sum_{q\in\mathcal Q_S}
       |\mathcal R_{\rm ref}(q)\setminus C_\alpha(q)|.
 \label{eq:6-reference-coverage}
\end{equation}
For valid fixed-size output sets, let $r_{\mathcal Q_S}$ denote GT recall
on these same queries. Each missing reference ID removes at most one GT hit, so
\begin{equation}
 r_{\mathcal Q_S}(\alpha)\geq r_{\mathcal Q_S}({\rm ref})
                         -\frac{D_S(\alpha)}{Sk}.
 \label{eq:6-relative-recall}
\end{equation}
This bound is sample-local and reference-relative. Agreement controls
additional loss relative to the reference; the reference's own quality
requires GT evaluation. \Cref{sec:7-recall-attribution} quantifies that
distinction with the reconstruction pilot.

\subsection{Sampling resolution and ordered testing}

Let $H_S$ count queries with any missing reference ID. Since $H_S\leq D_S$,
acceptance implies $H_S\leq1$. For a prespecified budget and i.i.d.
queries, the one-sided 95\% exact binomial upper limit after one event solves
\cite{nist_binomial}
\begin{equation}
 (1-u)^S+Su(1-u)^{S-1}=0.05.
 \label{eq:6-binomial-limit}
\end{equation}
The query-event limits are 0.03652 at $S=128$ and 0.009452 at 500.
Treating all $Sk$ result indicators as homogeneous independent Bernoulli
trials is stronger: the corresponding limits are 0.003701 and 0.0009484,
approximately $4.74/(Sk)$. Query i.i.d. alone does not justify dividing
the query-event limit by $k$; neither model certifies repeated subsamples
of this fixed pool.

For monotone disagreement and a fixed grid order, bisection selects the
accepted prefix. If its per-arm decisions are valid level tests, the
first-unsafe-arm argument gives fixed-sequence error control without a
Bonferroni factor \cite{learn_then_test}. The implemented one-slot rule
is a stability threshold; a GT-risk guarantee would require a valid test
for that target.

\paragraph{A finite-pool extension.}
Suppose every query has a fixed disagreement count $X_q(\alpha)$ under
the intended sample-size execution regime, unchanged by subset membership
or position. Let $N_{\rm cal}$ be the pool size and $n_j$ count queries with $X_q=j$.
Under uniform sampling without replacement, $D_S\leq1$ requires either
only zero-count queries or one single-count query \cite{nist_hypergeometric}:
\begin{equation}
 \Pr\{D_S(\alpha)\leq1\}
 =\frac{\binom{n_0}{S}+n_1\binom{n_0}{S-1}}{\binom{N_{\rm cal}}{S}}.
 \label{eq:6-finite-pool}
\end{equation}
Out-of-range binomial coefficients are zero. This is a conditional
acceptance model. Predicting selection risk additionally needs the
grid's acceptance structure and selected-arm quality map. The supplementary full-pool records reproduce every sampled curve by
summing the recorded $X_q$ over its actual sample IDs: 18,432 grid counts
match, including changes of execution path (\cref{sec:7-budget-trace}).
This validates the aggregate counts on those draws. Prediction for arbitrary
subsets still assumes stable per-query counts. Acceptance needs no labels;
predicting GT-loss crossings additionally uses the fixed-arm held-out
quality map.


\section{Budget Search, Nestedness and Tie Safety}
\label{sec:6-3}
\subsection{Bisection and complete enumeration}

Write $J=|A|$. The implementation computes the sample reference once,
then searches for the largest acceptable index in the descending grid:
\par\noindent\begin{minipage}{\linewidth}
\begin{enumerate}[leftmargin=*,itemsep=0.15em]
\item Set $l=0$, $h=J-1$. Index zero is the reference fallback.
\item While $l<h$, test $j=\lfloor(l+h+1)/2\rfloor$ by running
$A[j]$ and computing $D_S(A[j])$.
\item If $D_S(A[j])\leq\epsilon$, set $l=j$; otherwise set $h=j-1$.
\item Return $A[l]$. A caller applies it before processing reuse batches.
\end{enumerate}
\end{minipage}\par\medskip
The eight-point grid needs at most three complete candidate searches
in addition to the reference. Invocation, host-comparison complexity and
timing boundaries are specified in \cref{sec:a-budget-contract}.

Complete enumeration repeats the reference and tests all seven other
budgets on exactly the same sample, selecting the smallest accepted
multiplier. Including a fresh reference in both timers aligns their cost
boundaries. Enumeration evaluates the finite grid even if acceptance has
holes; bisection requires an accepted prefix to guarantee the same minimum.

\subsection{A sufficient condition for monotonicity}

\par\noindent\begin{minipage}{\linewidth}
\begin{proposition}[Nested survivors preserve reference membership]
\label{prop:6-monotonicity}
For one query, assume fixed routing and survivor sets
$C_\alpha\subseteq C_\beta\subseteq C_{\alpha_{\rm ref}}$ whenever
$\alpha\leq\beta\leq\alpha_{\rm ref}$. Assume at least $k$ valid
survivors, fixed refined scores and one deterministic total order for
final ties. Each result is the exact top-$k$ in that order. Then
\begin{equation}
 |\mathcal R_\alpha\cap\mathcal R_{\rm ref}|
 =|C_\alpha\cap\mathcal R_{\rm ref}|,
 \qquad D_S(\alpha)\geq D_S(\beta).
 \label{eq:6-monotonicity}
\end{equation}
The latter statement holds when these conditions hold for every sampled
query.
\end{proposition}
\begin{proof}
For $x\in\mathcal R_{\rm ref}\cap C_\alpha$, the common total order gives
\begin{equation*}
 |\{y\in C_\alpha:y\prec x\}|
 \leq|\{y\in C_{\alpha_{\rm ref}}:y\prec x\}|<k.
\end{equation*}
Hence $x\in\mathcal R_\alpha$. Together with
$\mathcal R_\alpha\subseteq C_\alpha$, this proves the intersection
identity in \cref{eq:6-monotonicity}. Nestedness gives
$|C_\alpha\cap\mathcal R_{\rm ref}|
 \leq|C_\beta\cap\mathcal R_{\rm ref}|$.
Summing missing reference counts over queries proves
$D_S(\alpha)\geq D_S(\beta)$.
\end{proof}
\end{minipage}\par\medskip

Exact primary top-$c$ under one fixed total key supplies the nesting
premise as $c$ grows. A fixed refined total key then makes acceptance
monotone, so acceptable indices form a prefix and bisection finds the
smallest accepted grid budget. Neither GT recall nor execution speed
appears in this theorem. Approximate refined scores can promote a false
neighbour as the pool grows, and runtime remains hardware dependent.

\begin{figure}[tbp]
\centering
% Synthetic sets only: all primary scores can tie; refined order is a<b<c.
\begingroup
\small\setlength{\tabcolsep}{7pt}
\begin{tabular}{@{}lcccc@{}}
\toprule
Case & $C_1$ & $C_2$ & $C_3$ & Reference overlap \\
\midrule
Nested & $\{b\}$ & $\{b,c\}$ & $\{a,b,c\}$ & $0,\ 0,\ 1$ \\
Non-nested & $\{a\}$ & $\{b,c\}$ & $\{a,b,c\}$ & $1,\ 0,\ 1$ \\
\bottomrule
\end{tabular}
\endgroup

\caption[Nested and non-nested candidate pools]{Constructed $k=1$ example with refined order $a\prec b\prec c$
and reference $\{a\}$. The nested case preserves reference membership.
With all primary scores tied, different tie choices can produce the
non-nested middle pool and break monotonicity. These are explanatory
sets, not measured GPU outputs.}
\label{fig:6-nesting}
\end{figure}

\subsection{A corollary for the evaluated grouped path}

The run-head source defaults to exact primary selection. Within one
calibration draw, $S$, routing and scan threads are fixed;
\texttt{scheduled\_groups(S)} is independent of $c$. Changing the budget
changes buffer capacity and admission work, while the grouped scorer
uses the same FP32 LUT entries and accumulation order. Refinement fixes
$B_r$ and 16 lanes, preserving each candidate's residual reduction;
final selection uses $(\refineddist,\mathrm{ID})$.

\par\noindent\begin{minipage}{\linewidth}
\begin{corollary}[Grouped calibration]
\label[corollary]{cor:6-grouped-calibration}
For the evaluated grouped configuration, under the representation and
shared-bound contracts of Chapters 4 and 5, primary survivors are nested
across budgets. Consequently \Cref{prop:6-monotonicity} applies and
bisection returns the same minimum accepted grid budget as enumeration.
\end{corollary}
\begin{proof}
\Cref{lem:5-buffer,prop:5-local-coverage,prop:5-shared-bound} give exact
primary top-$c$ under the common key $(\primarydist,\mathrm{position})$.
Its prefixes nest as $c$ grows. The fixed per-candidate refinement and
exact final total order supply the remaining premises of
\Cref{prop:6-monotonicity}.
\end{proof}
\end{minipage}\par\medskip

The corollary covers 2,176 of the 2,304 subsampled draws at
$S=32,64,128$. The remaining 128 use OpenAI-3072's distance-only
one-block selector at $S=128$. The supplementary trace records every
grid count: all 2,304 curves are monotone and all decisions match
enumeration (\cref{sec:7-budget-trace}).

\paragraph{Full-pool coverage.}
The twelve $S=500$ calls are separate: two Vogue cells use grouped
selection and ten use $\gscan=0$. The supplementary harness adds full-pool
enumeration: all twelve curves are monotone and all selected budgets
match the minimum accepted enumerated budget. This closes the finite-grid
empirical check. The ten distance-only cells still lie outside
\Cref{cor:6-grouped-calibration}; a common cutoff order, or no tie
straddling any grid cutoff including the reference cutoff, remains a
sufficient condition for a general nestedness guarantee. Matching the
held-out regime aids transfer without proving that broader condition.

\paragraph{A rank-based extension.}
Under the same nestedness and fixed-score conditions, let
$\operatorname{rank}_{\rm P}(x)$ be a reference answer's one-based rank under the
common primary key. Then
\begin{equation}
 D_S(\alpha)=\sum_{q\in\mathcal Q_S}\sum_{x\in\mathcal R_{\rm ref}(q)}
 \mathbf1\{\operatorname{rank}_{\rm P}(x)>c(\alpha)\}.
 \label{eq:6-rank-curve}
\end{equation}
Exporting these ranks would recover the grid curve from one reference
search. The current harness reruns each budget; rank extraction and
transfer remain prospective work under the same ordering conditions.


\section{Budget Baselines and Their Questions}
\label{sec:6-4}
\Cref{tab:6-controls} separates budget selection from its evaluation.
Fixed arms share the index, route and held-out queries; invocation and
timing details are in \cref{sec:a-budget-contract}.

\begin{table}[tbp]
\centering
\caption[Budget controls and their roles]{Budget controls and the questions they answer. Held-out GT is
used to assess every result, but only the diagnostic oracle uses it to
choose a budget. The reference is fixed at 200 for agreement-based arms.}
\label{tab:6-controls}
\footnotesize\setlength{\tabcolsep}{4pt}
\begin{tabularx}{\linewidth}{@{}>{\raggedright\arraybackslash}p{2.7cm}*{2}{>{\raggedright\arraybackslash}X}@{}}
\toprule
Control & Selection information & Question \\
\midrule
Fixed 100 & No calibration & Gain relative to the inherited default? \\
Fixed 64 and smaller grid arms & No calibration & Does a simple fixed choice remove the need for calibration? \\
Sampled bisection & $S$ unlabelled queries; slot agreement & How well does a low-cost sample choose a reusable budget? \\
Same-sample enumeration & Same $S$, reference and threshold; all grid arms & Did bisection miss the smallest accepted grid budget? \\
Full-pool calibration & All 500 calibration queries; same bisection rule & What changes with broader coverage and a stricter normalised allowance? \\
Held-out oracle & GT recall and timed QPS for all arms & Which measured arm is fastest under the diagnostic quality constraint? \\
\bottomrule
\end{tabularx}
\end{table}

Let $\widehat r(\alpha)$ and $\widehat Q(\alpha)$ be measured held-out
recall and QPS, and let $\widehat r^*=\max_{\alpha\in A}\widehat r(\alpha)$.
The oracle is
\begin{equation}
 \alpha_{\rm oracle}\in\operatorname*{arg\,max}_{\alpha\in A}
 \left\{\widehat Q(\alpha):
 \widehat r(\alpha)\geq\widehat r^*-\delta\right\},\qquad\delta=10^{-3}.
 \label{eq:6-oracle}
\end{equation}
The recall slack $\delta$ differs from the slot allowance $\epsilon$.
The oracle consumes held-out labels and timings: it is a finite-grid
diagnostic, not a deployable label-free decision rule.

Full-pool calibration remains label free and uses bisection, not a
GT-constrained scan of all budgets. Its larger sample also tightens
$\epsilon/(Sk)$. Same-sample enumeration instead holds sample coverage
and threshold fixed, isolating the search procedure. The oracle solves a
quality-constrained speed problem, whereas agreement minimises an accepted
multiplier. Even exact bisection and perfect calibration-pool coverage
need not select the oracle arm.

Fixed 64 tests whether a smaller universal default already captures the
benefit; fixed 200 supplies a high-budget control. Their observed quality
and selected-budget outcomes are in
\cref{sec:7-5,tab:7-budget-quality}.


\section{Hold-out Risk, Calibration Cost and Payback}
\label{sec:6-5}
\subsection{Choosing a calibration scale}

The full-pool working point is supported by the held-out results in
\cref{tab:7-budget-full}; sampled alternatives are evaluated in
\cref{fig:7-budget-risk}. The independent-slot model in \cref{sec:6-2}
provides a planning scale: at $k=10$, $4.74/(Sk)\approx10^{-3}$ suggests
$S$ of order 500, while 128 has coarser resolution. This strong model
concerns reference disagreement, not GT loss, and does not make the
empirical recommendation a prespecified quality guarantee.

The path and common-order conditions for FULL are addressed in
\cref{sec:6-3}.


\subsection{Risk and selected-arm performance}

Over $m$ calibration draws, let draw $r$ select $\widehat\alpha_r$.
Define held-out loss and threshold-crossing frequency by
\begin{equation}
 L_r=\widehat r^*-\widehat r(\widehat\alpha_r),\qquad
 \widehat p_\delta=\frac{1}{m}\sum_{r=1}^{m}\mathbf1\{L_r>\delta\}.
 \label{eq:6-risk}
\end{equation}
Mean loss, crossing frequency and tail size answer different questions.
\Cref{tab:7-budget-quality,fig:7-budget-risk} report the archived summaries;
the largest supplementary per-draw losses are reported in \cref{sec:7-budget-trace}.

Reported selected-arm rate is
$\overline Q=m^{-1}\sum_r\widehat Q(\widehat\alpha_r)$, a lookup from timed
fixed arms. It is distinct from the median budget's rate and from the
throughput of a calibrated query stream. Quality and speed must be read
together (\cref{tab:7-budget-cost}).

\subsection{Cost of reuse}

For $B_{\rm reuse}$ equal-sized subsequent batches, model
\begin{equation}
 T_{\rm policy}(B_{\rm reuse})=T_{\rm cal}+B_{\rm reuse}T_{\rm chosen},\qquad
 T_{\rm fixed}(B_{\rm reuse})=B_{\rm reuse}T_f.
 \label{eq:6-total-cost}
\end{equation}
When $T_{\rm chosen}<T_f$, the break-even count is
\begin{equation}
 B_{\rm reuse}^*=\frac{T_{\rm cal}}{T_f-T_{\rm chosen}}.
 \label{eq:6-payback}
\end{equation}
Positive saving requires $T_{\rm chosen}<T_f$ and
$B_{\rm reuse}>B_{\rm reuse}^*$. Batch time is $T=1000n/Q$ milliseconds.
When choices vary, $1/\mathbb E[Q]\leq\mathbb E[1/Q]$, so inverse mean
QPS understates mean time. Timing boundaries are in
\cref{sec:a-budget-contract}; \cref{sec:7-stream-cost} tests direct
calibration-plus-reuse cost on four cells with quality alongside time.
Campaign estimates are in \cref{tab:7-budget-cost}.

Recalibration every $H$ batches adds $T_{\rm cal}/H$ per batch.
Drift detection and automatic recalibration remain future work.


\section{Policy Failure Modes and Applicability}
\label{sec:6-6}
Reuse requires a stable index, $k$, routing and query population;
batch-dependent scheduling may change ties and arithmetic. Valid, unique
IDs and common ordering or matched enumeration underpin the decision
(\cref{sec:6-3}); held-out quality, tail loss and amortised cost determine
its utility (\cref{sec:6-5}). The harness does not detect workload drift.


The observed family favours FULL for quality coverage and $S=128$ for
lower calibration cost. The caller explicitly installs $\alpha$.
Chapter 7 evaluates fixed-budget/oracle controls and four direct reuse
measurements.

The reported budget campaign and its trace rerun use v69 arithmetic,
including TF32 on the OpenAI datasets. The separate FP32 sensitivity
measures a search frontier; it does not recalibrate or validate these
budget choices under FP32.

